\ifdefined\gkver\else\def\gkver{student}\fi
\documentclass[\gkver]{gaokaozhenti}
\begin{document}

\chapter{2001年全国旧课程}

\begin{enumerate}
\item
%% number: 1
%% paperName: 2001年全国高考旧课程第 1 题 4 分
%% typeId: 2
%% type: 多选题
%% chapter: 原子和原子核
%% point: 人工核反应
%% method: 无
%% score: 4
%% degree: 650
%% duplicateId: 0
%% body:
在下列的四个方程中，$x_{1}$、$x_{2}$、$x_{3}$ 和 $x_{4}$ 各代表某种粒子

① $\ce{^{235}_{92}U}+\ce{^{1}_{0}n}\to\ce{^{95}_{38}Sr}+\ce{^{138}_{54}Xe}+3x_{1}$　　② $\ce{^{2}_{1}H}+x_{2}\to\ce{^{3}_{2}He}+\ce{^{1}_{0}n}$

③ $\ce{^{238}_{92}U}\to\ce{^{234}_{90}Th}+x_{3}$　　④ $\ce{^{24}_{12}Mg}+\ce{^{4}_{2}He}\to\ce{^{27}_{13}Al}+x_{4}$

以下判断正确的是\xzanswer{AC}

\fourchoices[answer=AC]
{$x_{1}$ 是中子}
{$x_{2}$ 是质子}
{$x_{3}$ 是 $\alpha$ 粒子}
{$x_{4}$ 是氘核}

%% answer: AC
%% memo:
\memoanswer{
首先由核电荷数守恒算出 $x_{1}$、$x_{2}$、$x_{3}$、$x_{4}$ 的核电荷数分别为 0、1、2、1，然后再由质量数守恒确定 $x_{1}$ 为中子，$x_{2}$ 为氘核，$x_{3}$ 为 $\alpha$ 粒子，$x_{4}$ 为质子，故 AC 正确。
}

\item
%% number: 2
%% paperName: 2001年全国高考旧课程第 2 题 4 分
%% typeId: 2
%% type: 多选题
%% chapter: 交流电
%% point: 变压器
%% method: 无
%% score: 4
%% degree: 650
%% duplicateId: 0
%% body:
一个理想变压器，原线圈和副线圈的匝数分别是 $n_{1}$ 和 $n_{2}$，正常工作时输入和输出的电压、电流、功率分别为 $U_{1}$ 和 $U_{2}$、$I_{1}$ 和 $I_{2}$、$P_{1}$ 和 $P_{2}$，已知 $n_{1} > n_{2}$，则\xzanswer{BC}

\fourchoices[answer=BC]
{$U_{1} > U_{2}$，$P_{1} < P_{2}$}
{$P_{1} = P_{2}$，$I_{1} < I_{2}$}
{$I_{1} < I_{2}$，$U_{1} > U_{2}$}
{$P_{1} > P_{2}$，$I_{1} > I_{2}$}

%% answer: BC
%% memo:
\memoanswer{
由理想变压器的原理得 $\frac{U_{1}}{U_{2}}=\frac{n_{1}}{n_{2}}$，$P_{1}=P_{2}$，由题意知 $n_{1}>n_{2}$，则 $U_{1}>U_{2}$，又因 $I_{1}U_{1}=I_{2}U_{2}$，则 $I_{1}<I_{2}$，故 BC 正确。
}

\item
%% number: 3
%% paperName: 2001年全国高考旧课程第 3 题 4 分
%% typeId: 1
%% type: 单选题
%% chapter: 光学
%% point: 光子说
%% method: 无
%% score: 4
%% degree: 650
%% duplicateId: 0
%% body:
在 X 射线管中，由阴极发射的电子被加速后打到阳极，会产生包括 X 光在内的光子，其中光子能量的最大值等于电子的动能。已知阳极与阴极之间的电势差 $U$。普朗克常数 $h$，电子电量 $e$ 和光速 $c$，则可知该 X 射线管发出的 X 光的\xzanswer{D}

\onepicture[w=4.0cm]{figs/03.png}

\fourchoices[answer=D]
{最短波长为 $\frac{c}{eUh}$}
{最长波长为 $\frac{hc}{eU}$}
{最小频率为 $\frac{eU}{h}$}
{最大频率为 $\frac{eU}{h}$}

%% answer: D
%% memo:
\memoanswer{
由动能定理知加速电场对电子所做的功等于电子动能的增量，由题意知光子的最大能量等于电子的动能，则有 $h\nu_{m}=eU$，故 X 光的最大频率 $\nu_{m}=\frac{eU}{h}$，故 D 正确；X 光的最小波长为 $\lambda=\frac{c}{\nu_{m}}=\frac{ch}{eU}$，故 A 错误；因光子的最小能量无法确定，所以光子的最小频率和最长波长无法确定，故 BC 错误。
}

\item
%% number: 4
%% paperName: 2001年全国高考旧课程第 4 题 4 分
%% typeId: 1
%% type: 单选题
%% chapter: 光学
%% point: 透镜
%% method: 无
%% score: 4
%% degree: 750
%% duplicateId: 0
%% body:
如图所示，p 字形发光物经透镜 L 在毛玻璃光屏 M 上成一实像，观察者处于 E 处，他看到屏 M 上的像的形状为\xzanswer{C}

\onepicture[w=5.5cm]{figs/04.png}

\fourchoices[answer=C]
{q}
{p}
{d}
{b}

%% answer: C
%% memo:
\memoanswer{
由透镜成像原理知，透镜成实像左右、上下都倒置，故 C 正确。
}

\item
%% number: 5
%% paperName: 2001年全国高考旧课程第 5 题 4 分
%% typeId: 1
%% type: 单选题
%% chapter: 电磁感应
%% point: 电磁感应受力分析
%% method: 无
%% score: 4
%% degree: 550
%% duplicateId: 0
%% body:
如图所示，虚线框 $abcd$ 内为一矩形匀强磁场区域，$ab=2bc$，磁场方向垂直于纸面，线框 $a'b'c'd'$ 是一正方形导线框，$a'b'$ 边与 $ab$ 边平行。若将导线框匀速地拉离磁场区域，以 $W_{1}$ 表示沿平行于 $ab$ 的方向拉出过程中外力所做的功，$W_{2}$ 表示以同样的速率沿平行于 $bc$ 的方向拉出过程中外力所做的功，则\xzanswer{B}

\onepicture[w=5.0cm]{figs/05.png}

\fourchoices[answer=B]
{$W_{1} = W_{2}$}
{$W_{2} = 2W_{1}$}
{$W_{1} = 2W_{2}$}
{$W_{2} = 4W_{1}$}

%% answer: B
%% memo:
\memoanswer{
设导线框运动速率为 $v$、电阻为 $R$。由于导线框被匀速拉离磁场区域，所以外力所做的功等于安培力做的功，则有

$W_{1}=F_{\text{安}1}\cdot\overline{ab}=B\frac{B\overline{bc}v}{R}\cdot\overline{bc}\cdot\overline{ab}=\frac{B^{2}v\overline{bc}^{2}\overline{ab}}{R}$；

同理

$W_{2}=F_{\text{安}2}\cdot\overline{bc}=B\frac{B\overline{ab}v}{R}\cdot\overline{ab}\cdot\overline{bc}=\frac{B^{2}v\overline{ab}^{2}\overline{bc}}{R}$。

因为 $\overline{ab}=2\overline{bc}$，所以 $W_{2}=2W_{1}$，故 B 正确。
}

\item
%% number: 6
%% paperName: 2001年全国高考旧课程第 6 题 4 分
%% typeId: 2
%% type: 多选题
%% chapter: 原子和原子核
%% point: 玻尔模型
%% method: 无
%% score: 4
%% degree: 650
%% duplicateId: 0
%% body:
按照玻尔理论，下列关于氢原子的论述正确的是\xzanswer{AB}

\fourchoices[answer=AB]
{第 $m$ 个定态和第 $n$ 个定态的轨道半径 $r_{m}$ 和 $r_{n}$ 之比为 $r_{m}:r_{n} = m^{2}:n^{2}$}
{第 $m$ 个定态和第 $n$ 个定态的能量 $E_{m}$ 和 $E_{n}$ 之比为 $E_{m}:E_{n} = n^{2}:m^{2}$}
{电子沿某一轨道绕核运动，若其圆周运动的频率是 $\nu$，则其发光频率也是 $\nu$}
{若氢原子处于能量为 $E$ 的定态，则其发光频率为 $\nu = \frac{E}{h}$}

%% answer: AB
%% memo:
\memoanswer{
由氢原子电子轨道半径公式 $r_{n}=n^{2}r_{1}$ 知，轨道半径与量子数的平方成正比，所以 $r_{m}:r_{n}=m^{2}:n^{2}$，故 A 正确；由氢原子的能量公式 $E_{n}=\frac{1}{n^{2}}E_{1}$ 知，氢原子的能量与量子数的平方成反比，所以 $E_{m}:E_{n}=n^{2}:m^{2}$，故 B 正确；由玻尔的原子理论知，只有核外电子发生能级跃迁时才可能发射某种频率的光，故 CD 错误。
}

\item
%% number: 7
%% paperName: 2001年全国高考旧课程第 7 题 4 分
%% typeId: 2
%% type: 多选题
%% chapter: 电场
%% point: 带电粒子的轨迹类问题
%% method: 无
%% score: 4
%% degree: 550
%% duplicateId: 0
%% body:
如图，虚线 $a$、$b$ 和 $c$ 是某静电场中的等势面，它们的电势分别为 $\varphi_{a}$、$\varphi_{b}$ 和 $\varphi_{c}$，$\varphi_{a} > \varphi_{b} > \varphi_{c}$。一带正电的粒子射入电场中，其运动轨迹如实线 KLMN 所示，由图可知\xzanswer{AC}

\onepicture[w=4.6cm]{\includesvg{figs/07}}

\fourchoices[answer=AC]
{粒子从 K 到 L 的过程中，电场力做负功}
{粒子从 L 到 M 的过程中，电场力做负功}
{粒子从 K 到 L 的过程中，静电势能增加}
{粒子从 L 到 M 的过程中，动能减少}

%% answer: AC
%% memo:
\memoanswer{
由题图 $\varphi_{a}$、$\varphi_{b}$、$\varphi_{c}$ 的关系可知，静电场是正的点电荷产生的电场，当带正电粒子从 K 到 L 的过程中，粒子受到库仑力的作用，它克服电场力做功，其动能减少，电势能增加，故 AC 正确；当带正电粒子从 L 到 M 的过程中，电场力做正功，粒子的动能增加，电势能减少，故 BD 错误。
}

\item
%% number: 8
%% paperName: 2001年全国高考旧课程第 8 题 4 分
%% typeId: 1
%% type: 单选题
%% chapter: 运动和力
%% point: 牛顿第二定律
%% method: 无
%% score: 4
%% degree: 550
%% duplicateId: 0
%% body:
惯性制导系统已广泛应用于弹道式导弹工程中，这个系统的重要元件之一是加速度计，加速度计的构造原理的示意图如图所示：沿导弹长度方向安装的固定光滑杆上套一质量为 $m$ 的滑块，滑块两侧分别与劲度系数均为 $k$ 的弹簧相连；两弹簧的另一端与固定端相连，滑块原来静止，弹簧处于自然长度，滑块上有指针，可通过标尺测出滑块的位移，然后通过控制系统进行制导，设某段时间内导弹沿水平方向运动，指针向左偏离 O 点的距离为 $x$，则这段时间内导弹的加速度\xzanswer{D}

\onepicture[w=6.0cm]{\includesvg{figs/08}}

\fourchoices[answer=D]
{方向向左，大小为 $\frac{kx}{m}$}
{方向向右，大小为 $\frac{kx}{m}$}
{方向向左，大小为 $\frac{2kx}{m}$}
{方向向右，大小为 $\frac{2kx}{m}$}

%% answer: D
%% memo:
\memoanswer{
滑块随导弹一起做加速运动，向左偏离 O 点距离为 $x$，使左侧弹簧被压缩，右侧弹簧被拉长，则滑块所受合力为 $2kx$，方向向右，由牛顿第二定律得 $2kx=ma$，滑块加速度大小为 $a=\frac{2kx}{m}$，故 D 正确。
}

\item
%% number: 9
%% paperName: 2001年全国高考旧课程第 9 题 4 分
%% typeId: 2
%% type: 多选题
%% chapter: 机械振动
%% point: 单摆
%% method: 无
%% score: 4
%% degree: 450
%% duplicateId: 0
%% body:
细长轻绳下端栓一小球构成单摆，在悬挂点正下方 $\frac{1}{2}$ 摆长处有一个能挡住摆线的钉子 A，如图所示，现将单摆向左方拉开一个小角度，然后无初速度地释放，对于以后的运动，下列说法正确的是\xzanswer{AB}

\onepicture[h=5.0cm]{\includesvg{figs/09}}

\fourchoices[answer=AB]
{摆球往返运动一次的周期比无钉子时的单摆周期小}
{摆球在左、右两侧上升的最大高度一样}
{摆球在平衡位置左右两侧走过的最大弧长相等}
{摆球在平衡位置右侧的最大摆角是左侧的两倍}

%% answer: AB
%% memo:
\memoanswer{
由单摆的周期公式 $T=2\pi\sqrt{\frac{l}{g}}$ 知，有钉子时，$l'=\frac{l}{2}$，故周期变小，故 A 正确；由机械能守恒 $\frac{1}{2}mv^{2}=mgh$ 知，摆球上升的高度 $h$ 与摆角无关，故 B 正确；如图所示，假如没有钉子 A，则摆球可以摆到 E（E 与 B 关于 OAC 竖直线对称），从图中可以看出，弧长 BC 大于弧长 CD，故 C 错误；由图知，$\angle CAD=2\angle AOD$，$\angle BOC>\angle AOD$，所以 $\angle CAD<2\angle BOC$，故 D 错误。

\onepicture[h=6.0cm]{figs/09_an.png}
}

\item
%% number: 10
%% paperName: 2001年全国高考旧课程第 10 题 4 分
%% typeId: 2
%% type: 多选题
%% chapter: 机械波
%% point: 波的图象
%% method: 无
%% score: 4
%% degree: 450
%% duplicateId: 0
%% body:
如图所示，在平面 $xy$ 内有一沿水平轴 $x$ 正向传播的简谐横波，波速为 $3.0\Ums$，频率为 $2.5\UHz$，振幅为 $8.0\times10^{-2}\Um$。已知 $t=0$ 时刻 P 质点的位移为 $y=4.0\times10^{-2}\Um$，速度沿 $y$ 轴正向。Q 点在 P 点右方 $9.0\times10^{-1}\Um$ 处，对于 Q 点的质元来说\xzanswer{BC}

\onepicture[w=7.0cm]{\includesvg{figs/10}}

\fourchoices[answer=BC]
{在 $t=0$ 时，位移为 $y=-4.0\times10^{-2}\Um$}
{在 $t=0$ 时，速度沿 $y$ 轴负方向}
{在 $t=0.1\Us$ 时，位移为 $y=-4.0\times10^{-2}\Um$}
{在 $t=0.1\Us$ 时，速度沿 $y$ 轴正方向}

%% answer: BC
%% memo:
\memoanswer{
首先画出一列波形图，如图中实线所示，在图中标出 $y=4.0\times10^{-2}\Um$ 并且速度方向沿 $y$ 轴正方向的 P 点，由 $v=\lambda f$ 得 $\lambda=\frac{v}{f}=\frac{3.0}{2.5}\Um=1.2\Um$，Q 点距 P 点 $\Delta x=0.9\Um=\frac{3}{4}\lambda$，从而找出 Q 点标在图中，从图中 Q 点的位置可以看出，故 A 错误 B 正确。由 $T=\frac{1}{f}=0.4\Us$ 得 $t=0.1\Us=\frac{1}{4}T$，所以波在 $t=0.1\Us$ 时向右传了 $\frac{1}{4}\lambda$，据此作出 $t=0.1\Us$ 时的波形图，如图中虚线所示，从图中可以看出，此时 Q 点的位移 $y=-4\times10^{-2}\Um$，速度方向沿 $y$ 轴负方向，故 C 正确 D 错误。

\onepicture[w=8.0cm]{figs/10_an.png}
}

\gknewpage

\item
%% number: 11
%% paperName: 2001年全国高考旧课程第 11 题 5 分
%% typeId: 3
%% type: 填空题
%% chapter: 直线运动
%% point: 匀速直线运动
%% method: 无
%% score: 5
%% degree: 650
%% duplicateId: 0
%% body:
某测量员是这样利用回声测距离的：他站在两平行峭壁间某一位置鸣枪，经过 $1.00\Us$ 第一次听到回声，又经过 $0.50\Us$ 再次听到回声。已知声速为 $340\Ums$，则两峭壁间的距离为\tkanswer[2.6cm]{$425\Um$}。

%% answer: 
%% memo:
\memoanswer{
由声波传播特性知，$L=340\times\frac{1.00}{2}+340\times\frac{1.50}{2}=425\Um$。
}

\item
%% number: 12
%% paperName: 2001年全国高考旧课程第 12 题 5 分
%% typeId: 3
%% type: 填空题
%% chapter: 力物体的平衡
%% point: 多力平衡
%% method: 无
%% score: 5
%% degree: 850
%% duplicateId: 0
%% body:
如图所示，质量为 $m$，横截面为直角三角形的物块 $ABC$，$\angle ABC=\alpha$。$AB$ 边靠在竖直墙面上，$F$ 是垂直于斜面 $BC$ 的推力。现物块静止不动，则摩擦力的大小为\tkanswer[3.0cm]{$mg+F\sin\alpha$}。

\onepicture[h=3.6cm, align=right]{\includesvg{figs/12}}

%% answer: 
%% memo:
\memoanswer{
对物块受力分析，作出受力分析图，由平衡条件得，摩擦力 $f=mg+F\sin\alpha$。

\onepicture[h=3.8cm]{figs/12_an.png}
}

\item
%% number: 13
%% paperName: 2001年全国高考旧课程第 13 题 5 分
%% typeId: 3
%% type: 填空题
%% chapter: 电场
%% point: 带电粒子的平衡
%% method: 无
%% score: 5
%% degree: 350
%% duplicateId: 0
%% body:
如图所示，$q_{1}$、$q_{2}$、$q_{3}$ 分别表示在一条直线上的三个点电荷，已知 $q_{1}$ 与 $q_{2}$ 之间的距离为 $l_{1}$，$q_{2}$ 与 $q_{3}$ 之间的距离为 $l_{2}$，且每个电荷都处于平衡状态

\begin{enumerate}
	\item
	如 $q_{2}$ 为正电荷，则 $q_{1}$ 为\_\_\_\_\_\_电荷，$q_{3}$ 为\_\_\_\_\_\_电荷。
	\item
	$q_{1}$、$q_{2}$、$q_{3}$ 三者电量大小之比是\_\_\_\_\_\_：\_\_\_\_\_\_：\_\_\_\_\_\_。
\end{enumerate}

\onepicture[w=7.0cm, align=right]{\includesvg{figs/13}}

%% answer: 
\jdanswer{
\begin{enumerate}
	\item
	负，负
	\item
	$\left(\frac{l_{1}+l_{2}}{l_{2}}\right)^{2}:1:\left(\frac{l_{1}+l_{2}}{l_{1}}\right)^{2}$
\end{enumerate}
}
%% memo:
\memoanswer{
（1）若 $q_{2}$ 为正电荷，且每个电荷都处于平衡状态，则 $q_{1}$、$q_{3}$ 均要带负电荷才能满足要求。

（2）由于三个点电荷都处于平衡状态，因此三个点电荷所在处的合场强均为零，即有

$k\frac{q_{2}}{l_{1}^{2}}=k\frac{q_{3}}{(l_{1}+l_{2})^{2}}$，$k\frac{q_{1}}{l_{1}^{2}}=k\frac{q_{3}}{l_{2}^{2}}$，$k\frac{q_{1}}{(l_{1}+l_{2})^{2}}=k\frac{q_{2}}{l_{2}^{2}}$，

用 $q_{2}$ 分别表示 $q_{1}$、$q_{3}$，可得

$q_{3}=q_{2}\left(\frac{l_{1}+l_{2}}{l_{1}}\right)^{2}$，$q_{1}=q_{2}\left(\frac{l_{1}+l_{2}}{l_{2}}\right)^{2}$，

所以 $q_{1}:q_{2}:q_{3}=\left(\frac{l_{1}+l_{2}}{l_{2}}\right)^{2}:1:\left(\frac{l_{1}+l_{2}}{l_{1}}\right)^{2}$。
}

\item
%% number: 14
%% paperName: 2001年全国高考旧课程第 14 题 5 分
%% typeId: 4
%% type: 实验题
%% chapter: 光学
%% point: 光的衍射
%% method: 无
%% score: 5
%% degree: 750
%% duplicateId: 0
%% body:
某同学以线状白炽灯为光源，利用游标卡尺两脚间形成的狭缝观察光的衍射现象后，总结出以下几点：

（A）若狭缝与灯丝平行，衍射条纹与狭缝平行

（B）若狭缝与灯丝垂直，衍射条纹与狭缝垂直

（C）衍射条纹的疏密程度与狭缝宽度有关

（D）衍射条纹的间距与光的波长有关

以上几点中，你认为正确的是\tkanswer[2.2cm]{ACD}。

%% answer: 
%% memo:
\memoanswer{
用狭缝观察线状光源时，无论狭缝与灯丝平行还是垂直，衍射条纹总与狭缝平行，故 A 正确，B 错误；衍射条纹的疏密程度与狭缝的宽度有关，宽度越窄衍射现象越明显，衍射条纹越疏，故 C 正确；衍射条纹的间距与光的波长有关，波长越长，越容易发生衍射，衍射条纹间距越大，故 D 正确。
}

\item
%% number: 15
%% paperName: 2001年全国高考旧课程第 15 题 6 分
%% typeId: 4
%% type: 实验题
%% chapter: 直线运动
%% point: 打点计时器公式
%% method: 无
%% score: 6
%% degree: 650
%% duplicateId: 0
%% body:
一打点计时器固定在斜面上某处，一小车拖着穿过打点计时器的纸带从斜面上滑下，如图\ref{2001全国旧课程15a}所示，图\ref{2001全国旧课程15b}是打出的纸带的一段。

\begin{enumerate}
	\item
	已知打点计时器使用的交流电频率为 $50\UHz$，利用图\ref{2001全国旧课程15b}给出的数据可求出小车下滑的加速度 $a = $\_\_\_\_\_\_\_\_。
	\item
	为了求出小车在下滑过程中所受的阻力，还需测量的物理量有\_\_\_\_\_\_\_\_\_。用测得的量及加速度 $a$ 表示阻力的计算式为 $f = $\_\_\_\_\_\_\_\_\_\_。
\end{enumerate}

\twopicture[numstyle=chinese, wA=5.0cm, wB=8.5cm, align=right, labelA=2001全国旧课程15a, labelB=2001全国旧课程15b]
{figs/15a.png}
{figs/15b.png}

%% answer: 
\jdanswer{
\begin{enumerate}
	\item
	$4.00\Umsq$（$3.90\sim4.10\Umsq$ 之间都正确）
	\item
	小车质量 $m$；斜面上任意两点间距离 $l$ 及这两点的高度差 $h$；$f=\frac{mgh}{l}-ma$
\end{enumerate}
}
%% memo:
\memoanswer{
（1）用公式 $a=\frac{(s_{9}+s_{8}+s_{7}+s_{6})-(s_{5}+s_{4}+s_{3}+s_{2})}{32T^{2}}$ 求出加速度，$T=0.02\Us$，代入数据得 $a\approx4.00\Umsq$。

（2）由 $mg\sin\alpha-f=ma$ 和 $\sin\alpha=\frac{h}{l}$，解得阻力 $f=mg\frac{h}{l}-ma$。
}

\item
%% number: 16
%% paperName: 2001年全国高考旧课程第 16 题 9 分
%% typeId: 4
%% type: 实验题
%% chapter: 电路
%% point: 电路实验
%% method: 无
%% score: 9
%% degree: 450
%% duplicateId: 0
%% body:
图\ref{2001全国旧课程16a}中 E 为电源，其电动势为 $E$，$R_{1}$ 为滑线变阻器，$R_{2}$ 为电阻箱，A 为电流表，用此电路，经以下步骤可近似得 A 的内阻 $R_{A}$：①闭合 $K_{1}$，断开 $K_{2}$，调节 $R_{1}$ 使电流表读数等于其量程 $I_{0}$；②保持 $R_{1}$ 不变，闭合 $K_{2}$，调节 $R_{2}$，使电流表读数等于 $\frac{I_{0}}{2}$，然后读出 $R_{2}$ 的值，取 $R_{A}\approx R_{2}$。

\begin{enumerate}
	\item
	按图\ref{2001全国旧课程16a}所示电路在图\ref{2001全国旧课程16b}中所给的实物图中画出连接导线。
	\item
	真实值和测量值之差除以真实值叫测量结果的相对误差，即 $\frac{R_{A}-R_{2}}{R_{A}}$，试导出它与电动势 $E$，电流表量程 $I_{0}$ 及电流表内阻 $R_{A}$ 的关系式。
	\item
	若 $I_{0}=10\UmA$，真实值 $R_{A}$ 约为 $30\UO$，要想使测量结果相对误差不大于 $5\%$，电源电动势最小应为多少伏？
\end{enumerate}

\twopicture[numstyle=chinese, wA=6.0cm, wB=6.5cm, align=right, labelA=2001全国旧课程16a, labelB=2001全国旧课程16b]
{figs/16a.png}
{figs/16b.png}

%% answer: 
\jdanswer{
\begin{enumerate}
	\item
	连线如图

	\onepicture[w=5.0cm]{figs/16_an.png}

	\item
	$\frac{R_{A}-R_{2}}{R_{A}}=\frac{I_{0}}{E}R_{A}$
	\item
	$6\UV$
\end{enumerate}
}
%% memo:
\memoanswer{
（2）推导过程如下：

由步骤①得 $\frac{E}{R_{1}+R_{A}}=I_{0}$ ①

由步骤②得 $\frac{E}{R_{1}+\frac{R_{A}R_{2}}{R_{A}+R_{2}}}\cdot\frac{R_{2}}{R_{A}+R_{2}}=\frac{I_{0}}{2}$ ②

②式乘以 2 等于①式，解得 $R_{A}=\frac{R_{1}R_{2}}{R_{1}-R_{2}}$ ③，

从而得 $R_{2}=\frac{R_{A}R_{1}}{R_{A}+R_{1}}$ ④，

因为 $\frac{R_{A}-R_{2}}{R_{A}}=1-\frac{R_{2}}{R_{A}}$，将④代入得

$\frac{R_{A}-R_{2}}{R_{A}}=1-\frac{R_{1}}{R_{A}+R_{1}}=\frac{R_{A}}{R_{A}+R_{1}}$ ⑤，

又据①得 $R_{1}=\frac{E}{I_{0}}-R_{A}$ ⑥，

将⑥代入⑤得 $\frac{R_{A}-R_{2}}{R_{A}}=\frac{I_{0}}{E}R_{A}$ ⑦。

（3）根据 $\frac{I_{0}}{E}R_{A}=0.05$，将 $I_{0}=10\UmA=0.010\UA$，$R_{A}=30\UO$ 代入解得 $E=6\UV$。
}

\item
%% number: 17
%% paperName: 2001年全国高考旧课程第 17 题 12 分
%% typeId: 6
%% type: 计算题
%% chapter: 运动和力
%% point: 动量守恒定律
%% method: 无
%% score: 12
%% degree: 650
%% duplicateId: 0
%% body:
质量为 $M$ 的小船以速度 $v_{0}$ 行驶，船上有两个质量皆为 $m$ 的小孩 a 和 b，分别静止站在船头和船尾，现小孩 a 沿水平方向以速率 $v$（相对于静止水面）向前跃入水中，然后小孩 b 沿水平方向以相同的速率 $v$（相对于静止水面）向后跃入水中，求小孩 b 跃出后小船的速度。

%% answer: 
\jdanswer{
$v' = \left(1 + \frac{2m}{M}\right)v_{0}$
}
%% memo:
\memoanswer{
设小孩 b 跃出后小船向前行驶的速度为 $v'$，根据动量守恒定律，有

$(M+2m)v_{0}=Mv'+mv-mv$ ①

解得 $v'=\left(1+\frac{2m}{M}\right)v_{0}$ ②

评分标准：本题 12 分。①式 8 分，②式 4 分。
}

\item
%% number: 18
%% paperName: 2001年全国高考旧课程第 18 题 12 分
%% typeId: 6
%% type: 计算题
%% chapter: 磁场
%% point: 带电粒子在磁场中的运动
%% method: 无
%% score: 12
%% degree: 550
%% duplicateId: 0
%% body:
如图所示，在 $y < 0$ 的区域内存在匀强磁场，磁场方向垂直于 $xy$ 平面并指向纸面外，磁感强度为 $B$，一带正电的粒子以速度 $v_{0}$ 从 O 点射入磁场，入射方向在 $xy$ 平面内，与 $x$ 轴正方向的夹角为 $\theta$，若粒子射出磁场的位置与 O 点的距离为 $l$，求该粒子的电量和质量之比 $\frac{q}{m}$。

\onepicture[h=4.4cm, align=right]{figs/18.png}

%% answer: 
\jdanswer{
$\frac{q}{m} = \frac{2v_{0}\sin\theta}{lB}$
}
%% memo:
\memoanswer{
带正电粒子射入磁场后，由于受到洛伦兹力的作用，粒子将沿图示的轨迹运动，从 A 点射出磁场，O、A 间的距离为 $l$，射出时速度的大小仍为 $v_{0}$，射出方向与 $x$ 轴的夹角仍为 $\theta$，由洛伦兹力公式和牛顿定律可得：$qv_{0}B=m\frac{v_{0}^{2}}{R}$，式中 $R$ 为圆轨道的半径。

解得 $R=\frac{mv_{0}}{qB}$ ①

圆轨道的圆心位于 OA 的中垂线上，由几何关系可得 $\frac{l}{2}=R\sin\theta$ ②

联立①、②两式，解得 $\frac{q}{m}=\frac{2v_{0}\sin\theta}{lB}$ ③

评分标准：本题 12 分。①式 4 分，②式 6 分，③式 2 分。

\onepicture[h=4.4cm]{figs/18_an.png}
}

\item
%% number: 19
%% paperName: 2001年全国高考旧课程第 19 题 12 分
%% typeId: 6
%% type: 计算题
%% chapter: 功和能
%% point: 能量守恒定律
%% method: 无
%% score: 12
%% degree: 450
%% duplicateId: 0
%% body:
“和平号”空间站已于今年 3 月 23 日成功地坠落在太平洋海域，坠落过程可简化为从一个近圆轨道（可近似看作圆轨道）开始，经过与大气摩擦，空间站的绝大部分经过升温、熔化，最后汽化而销毁，剩下的残片坠入大海。此过程中，空间站原来的机械能中，除一部分用于销毁和一部分被残片带走外，还有一部分能量 $E'$ 通过其他方式散失（不考虑坠落过程中化学反应的能量）

坠落开始时空间站的质量 $M = 1.17\times10^{5}\Ukg$；

轨道离地的高度为 $h = 146\Ukm$，地球半径 $R = 6.4\times10^{6}\Um$；

坠落范围内重力加速度可看作 $g = 10\Umsq$；

入海残片的质量 $m = 1.2\times10^{4}\Ukg$；

入海残片的温升高 $\Delta T = 3000\UK$；

入海残片的入海速度为声速 $v = 340\Ums$；

空间站材料每 1 千克升温 $1\UK$ 平均所需能量 $c = 1.0\times10^{3}\UJ$；

每销毁 1 千克材料平均所需能量 $\mu = 1.0\times10^{7}\UJ$。

\begin{enumerate}
	\item
	试导出以下各物理量的符号表示散失能量 $E'$ 的公式。
	\item
	算出 $E'$ 的数值。（结果保留两位有效数字）
\end{enumerate}

%% answer: 
\jdanswer{
\begin{enumerate}
	\item
	$E' = Mg\left(\frac{1}{2}R + \frac{3}{2}h\right) - (M-m)\mu - \frac{1}{2}mv^{2} - cm\Delta T$
	\item
	$E' = 2.9\times10^{12}\UJ$
\end{enumerate}
}
%% memo:
\memoanswer{
（1）根据题给条件，从近圆轨道到地面的空间中，重力加速度 $g=10\Umsq$。若以地面为重力势能零点，坠落过程开始时，空间站在近圆轨道的势能为

$E_{p}=Mgh$ ①

以 $v$ 表示空间站在近圆轨道上的速度，由牛顿定律可得

$Mg=M\frac{v^{2}}{r}$ ②

其中 $r$ 为轨道半径，以 $R_{\text{地}}$ 表示地球半径，则

$r=R+h$ ③

由②③式可得空间站在近圆轨道上的动能为

$E_{k}=\frac{1}{2}Mg(R+h)$ ④

由①④式得，在近圆轨道上空间站的机械能

$E=Mg\left(\frac{1}{2}R+\frac{3}{2}h\right)$ ⑤

在坠落过程中，用于销毁部分所需的能量为

$Q_{\text{汽}}=(M-m)\mu$ ⑥

用于残片升温所需的能量

$Q_{\text{残}}=cm\Delta T$ ⑦

残片的动能为 $E_{\text{残}}=\frac{1}{2}mv^{2}$ ⑧

以 $E'$ 表示其他方式散失的能量，则由能量守恒得

$E=Q_{\text{汽}}+E_{\text{残}}+Q_{\text{残}}+E'$ ⑨

由此得

$E'=Mg\left(\frac{1}{2}R+\frac{3}{2}h\right)-(M-m)\mu-\frac{1}{2}mv^{2}-cm\Delta T$ ⑩

（2）以题给数据代入得 $E'=2.9\times10^{12}\UJ$。

评分标准：本题 12 分。⑤式 4 分，⑥⑦⑧式各 1 分，⑨式 2 分，⑩式 1 分。
}

\item
%% number: 20
%% paperName: 2001年全国高考旧课程第 20 题 13 分
%% typeId: 6
%% type: 计算题
%% chapter: 电磁感应
%% point: 电磁感应中的匀变速
%% method: 无
%% score: 13
%% degree: 450
%% duplicateId: 0
%% body:
如图\ref{2001全国旧课程20a}所示，一对平行光滑轨道放置在水平面上，两轨道间距 $l = 0.20\Um$，电阻 $R = 1.0\UO$，有一导体杆静止放在轨道上，与两轨道垂直，杆及轨道的电阻可忽略不计，整个装置处于磁感强度 $B = 0.50\UT$ 的匀强磁场中，磁场方向垂直轨道面向下，现用一外力 $F$ 沿轨道方向拉杆，使之做匀加速运动，测得力 $F$ 与时间 $t$ 的关系如图\ref{2001全国旧课程20b}所示，求杆的质量 $m$ 和加速度 $a$。

\twopicture[numstyle=arabic, hA=3.4cm, hB=4.8cm, align=right, labelA=2001全国旧课程20a, labelB=2001全国旧课程20b]
{\includesvg{figs/20a}}
{\includesvg{figs/20b}}

%% answer: 
\jdanswer{
$a = 10\Umsq$，$m = 0.1\Ukg$
}
%% memo:
\memoanswer{
导体杆在轨道上做匀加速直线运动，用 $v$ 表示其速度，$t$ 表示时间，则有

$v=at$ ①

杆切割磁力线，将产生感应电动势

$E=Blv$ ②

在杆、轨道和电阻的闭合回路中产生电流

$I=\frac{E}{R}$ ③

杆受到的安培力为

$F_{\text{安}}=BIl$ ④

根据牛顿第二定律，有

$F-F_{\text{安}}=ma$ ⑤

联立以上各式，得

$F=\frac{B^{2}l^{2}a}{R}t+ma$ ⑥

由图线上取两点代入⑥式，可解得

$a=10\Umsq$，$m=0.1\Ukg$。
}

\item
%% number: 21
%% paperName: 2001年全国高考旧课程第 21 题 13 分
%% typeId: 6
%% type: 计算题
%% chapter: 气体性质
%% point: 查理定律
%% method: 无
%% score: 13
%% degree: 450
%% duplicateId: 0
%% body:
在一密封的啤酒瓶中，下方为溶有 $\ce{CO2}$ 的啤酒，上方为纯 $\ce{CO2}$ 气体，在 $20\UduC$ 时，溶解于啤酒中的 $\ce{CO2}$ 的质量为 $m_{A} = 1.050\times10^{-3}\Ukg$，上方气体状态 $\ce{CO2}$ 质量为 $m_{B} = 0.137\times10^{-3}\Ukg$，压强为 $p_{0} = 1$ 个大气压，当温度升高到 $40\UduC$ 时，啤酒中溶解的 $\ce{CO2}$ 的质量减少，变为 $m_{A}' = m_{A} - \Delta m$，瓶中气体 $\ce{CO2}$ 的压强上升到 $p_{1}$，已知 $\frac{m_{A}'}{m_{A}} = 0.60\times\frac{p_{1}}{p_{0}}$。啤酒的体积不因溶入 $\ce{CO2}$ 而变化，且不考虑容器体积和啤酒体积随温度变化，又知对同种气体，在体积不变的情况下 $\frac{p}{T}$ 与 $m$ 成正比，试计算 $p_{1}$ 等于多少标准大气压（结果保留两位有效数字）。

%% answer: 
\jdanswer{
$p_{1} = 1.6$ 标准大气压
}
%% memo:
\memoanswer{
在 $40\UduC$ 时，溶入啤酒的 $\ce{CO2}$ 的质量为

$m_{A}'=m_{A}-\Delta m$ ①

因质量守恒，气态 $\ce{CO2}$ 的质量为

$m_{B}'=m_{B}+\Delta m$ ②

由题设

$\frac{m_{A}'}{m_{A}}=0.60\times\frac{p_{1}}{p_{0}}$ ③

由于对同种气体，体积不变时，$\frac{p}{T}$ 与 $m$ 成正比，可得

$\frac{p_{1}}{p_{0}}=\frac{m_{B}'\times313}{m_{B}\times293}$ ④

由以上各式解得

$p_{1}=\frac{1+\frac{m_{A}}{m_{B}}}{0.6\frac{m_{A}}{m_{B}}+\frac{293}{313}}p_{0}$ ⑤

算得 $p_{1}=1.6\Uatm$ ⑥

评分标准：本题 13 分。②式 3 分，④式 4 分，⑤式 4 分，⑥式 2 分。
}

\item
%% number: 22
%% paperName: 2001年全国高考旧课程第 22 题 13 分
%% typeId: 6
%% type: 计算题
%% chapter: 功和能
%% point: 动能定理
%% method: 无
%% score: 13
%% degree: 250
%% duplicateId: 0
%% body:
一个圆柱形的竖直的井里存有一定量的水，井的侧面和底部是密闭的，在井中固定地插着一根两端开口的薄壁圆管，管和井共轴，管下端未触及井底，在圆管内有一不漏气的活塞，它可沿圆管上下滑动，开始时，管内外水面相齐，且活塞恰好触及水面，如图所示，现用卷扬机通过绳子对活塞施加一个向上的力 $F$，使活塞缓慢向上移动，已知管筒半径 $r = 0.100\Um$，井的半径 $R = 2r$，水的密度 $\rho = 1.00\times10^{3}\Ukgmc$，大气压强为 $p_{0} = 1.00\times10^{5}\UPa$，求活塞上升 $H = 9.00\Um$ 的过程中拉力 $F$ 所做的功。（井和管在水面以上及水面以下的部分足够长，不计活塞质量，不计摩擦，重力加速度 $g = 10\Umsq$）。

\onepicture[h=5.5cm, align=right]{figs/22.png}

%% answer: 
\jdanswer{
$1.65\times10^{4}\UJ$
}
%% memo:
\memoanswer{
从开始提升活塞至内外水面高度差为 $h_{0}=\frac{p_{0}}{\rho g}=10\Um$ 的过程中，活塞始终与管内液体接触。（再提升活塞时，活塞和水面之间将出现真空，另行讨论。）设活塞上升距离为 $h_{1}$，管外液面下降距离为 $h_{2}$，由几何关系得

$h_{0}=h_{1}+h_{2}$ ①

因流体体积不变，有

$h_{2}=h_{1}\frac{\pi r^{2}}{\pi R^{2}-\pi r^{2}}=\frac{1}{3}h_{1}$ ②

得 $h_{1}=\frac{3}{4}h_{0}=\frac{3}{4}\times10\Um=7.5\Um$ ③

题给 $H=9\Um>h_{1}$，由此可知确实有活塞下面是真空的一段过程。活塞移动距离从零到 $h_{1}$ 的过程中，对于水和活塞这个整体，其机械能的增量应等于除重力外其他力所做的功。因为始终无动能，所以机械能的增量也就等于重力势能增量，即：

$\Delta E=\rho(\pi r^{2}h_{1})g\frac{h_{0}}{2}$ ④

其他力有管内、外的大气压力和拉力 $F$。因为液体不可压缩，所以管内外大气压做的总功 $p_{0}\pi(R^{2}-r^{2})h_{2}-p_{0}\pi r^{2}h_{1}=0$，故外力做功就只是拉力 $F$ 做的功，由功能关系知 $W_{1}=\Delta E$ ⑤

即 $W_{1}=\rho(\pi r^{2})g\frac{3}{8}h_{0}^{2}=\frac{3}{8}\pi r^{2}\frac{p_{0}^{2}}{\rho g}=1.18\times10^{4}\UJ$ ⑥

活塞移动距离从 $h_{1}$ 到 $H$ 的过程中，液面不变，$F$ 是恒力，$F=\pi r^{2}p_{0}$，做的功为 $W_{2}=F(H-h_{1})=\pi r^{2}p_{0}(H-h_{1})=4.71\times10^{3}\UJ$ ⑦

所求拉力 $F$ 做的总功为 $W_{1}+W_{2}=1.65\times10^{4}\UJ$ ⑧

\onepicture[h=6.0cm]{figs/22_an.png}
}

\end{enumerate}

\end{document}
